%% LyX 2.4.4 created this file.  For more info, see https://www.lyx.org/.
%% Do not edit unless you really know what you are doing.
\documentclass[english,t]{beamer}
\usepackage{mathptmx}
\usepackage[T1]{fontenc}
\usepackage[latin9]{inputenc}
\usepackage{amssymb}
\usepackage{cancel}

\makeatletter

%%%%%%%%%%%%%%%%%%%%%%%%%%%%%% LyX specific LaTeX commands.
%% Because html converters don't know tabularnewline
\providecommand{\tabularnewline}{\\}

%%%%%%%%%%%%%%%%%%%%%%%%%%%%%% Textclass specific LaTeX commands.
% this default might be overridden by plain title style
\newcommand\makebeamertitle{\frame{\maketitle}}%
% (ERT) argument for the TOC
\AtBeginDocument{%
  \let\origtableofcontents=\tableofcontents
  \def\tableofcontents{\@ifnextchar[{\origtableofcontents}{\gobbletableofcontents}}
  \def\gobbletableofcontents#1{\origtableofcontents}
}

%%%%%%%%%%%%%%%%%%%%%%%%%%%%%% User specified LaTeX commands.
\usetheme{Warsaw}
% or ...

\setbeamercovered{transparent}
% or whatever (possibly just delete it)
\setbeamercovered{invisible} %makes overlays invisible

%gets rid of navigation symbols
\setbeamertemplate{navigation symbols}{}

%\usepackage{hyperref}
%\usepackage[usenames,dvipsnames,table]{xcolor} %adds additional color options
\definecolor{links}{HTML}{2A1B81}
%\hypersetup{colorlinks,linkcolor=blue,urlcolor=links}

\usepackage{multicol}

\usepackage{tikz}
\usetikzlibrary{calc}
\usetikzlibrary{plotmarks}
\usetikzlibrary{shapes}
\usepackage{pgfplots}
\pgfplotsset{compat=newest}

\usepackage{forest} %%for tree diagram

\usepackage{calc}
\usepackage[nomessages]{fp}

%\usepackage{advdate}    % Advancing/saving dates
\usepackage{datetime}   % Dates formatting
%\usepackage{datenumber} % Counters for dates

\usepackage[super]{nth}

\usepackage{etex}

\usepackage{fancybox}

%%table colors
%\usepackage[table]{xcolor}
\usepackage{colortbl}

\usepackage[outline]{contour}  %allows text halos
\contourlength{1pt} %sets halo width
%%example: \node at (0.75,0.5) {\contour{white}{\Large Text!}};
%%example:  \node [fill=white,inner sep=1pt] at (2.25,0.5) {\Large 			Text!};
%%example:  \node [fill=white,rounded corners=2pt,inner sep=1pt] at 			(3.75,0.5) {\Large Text!};

%%%%%%%%%%%%%%%%%%%%%%%%%%%
\def\formatdateny{\csname noyear\languagename\expandafter\endcsname\formatdate}
\def\noyearenglish#1, \the\@year{#1}
\let\noyearamerican\noyearenglish
\let\noyearbritish\noyearenglish

%%allows uncover with forest diagrams
\tikzset{
    invisible/.style={opacity=0,text opacity=0},
    visible on/.style={alt=#1{}{invisible}},
    alt/.code args={<#1>#2#3}{%
      \alt<#1>{\pgfkeysalso{#2}}{\pgfkeysalso{#3}} % \pgfkeysalso doesn't change the path
    },
}
\forestset{
  visible on/.style={
    for tree={
      /tikz/visible on={#1},
      edge={/tikz/visible on={#1}}}}}

%%%redefines column alignment to match vertically
%\define@key{beamerframe}{t}[true]{% top
  %\beamer@frametopskip=.2cm plus .5\paperheight\relax%
  %\beamer@framebottomskip=0pt plus 1fill\relax%
 % \beamer@frametopskipautobreak=\beamer@frametopskip\relax%
  %\beamer@framebottomskipautobreak=\beamer@framebottomskip\relax%
  %\def\beamer@initfirstlineunskip{}%
%}

%%provides pdf with 4 slides per page; don't forget to add 'handout' to remove overlays
%\usepackage{pgfpages}
%\pgfpagesuselayout{4 on 1}[a4paper, landscape, border shrink=7mm]
%\pgfpageslogicalpageoptions{1}{border code=\pgfusepath{stroke}}
%\pgfpageslogicalpageoptions{2}{border code=\pgfusepath{stroke}}
%\pgfpageslogicalpageoptions{3}{border code=\pgfusepath{stroke}}
%\pgfpageslogicalpageoptions{4}{border code=\pgfusepath{stroke}}
%%%Don't forget to add 'handout'

\makeatother

\usepackage{babel}
\begin{document}
\newcommand{\xmax}{2000}
\newcommand{\xmin}{0}
\newcommand{\xstep}{0} %must be xmin+n
\newcommand{\ymax}{500}
\newcommand{\ymin}{0}
\newcommand{\ystep}{0} %must be ymin+n

\newcounter{countEx} \setcounter{countEx}{0}
\newcounter{countProb} \setcounter{countProb}{0}
\resetcounteronoverlays{countEx} \resetcounteronoverlays{countProb}

\newcommand{\ex}[3]{
	\addtocounter{countEx}{1}
	\only<#1-#2>{\begin{exampleblock} {Example \chNum.\secNum.\arabic{countEx}.} #3	\end{exampleblock}}
}

\newcommand{\prob}[4]{
	\addtocounter{countProb}{1}
	\only<#1-#2>{\begin{exampleblock} {Problem  \chNum.\secNum.\arabic{countProb}.\,#3} #4	\end{exampleblock}}
}

\newcommand{\yline}[4]{
	(#2 - #4)/(#1 - #3)*(x - #1) + #2
}

\beamerdefaultoverlayspecification{<+->}

%%%%Chapter and Section numbers%%%%%
\newcommand{\chNum}{1}
\newcommand{\secNum}{4}
\newcommand{\secTitle}{More on Conditionals}
\date{\formatdate{22}{9}{2014}}

\title[Section \chNum.\secNum: \secTitle]{Section \chNum.\secNum: \secTitle}

\subtitle{Finite Mathematics~~MTH 107~~Fall 2014}

\author{Dr.\,James Ashe}

\titlegraphic{\includegraphics[height=2.75cm]{/home/jim/JamesAshe/MHU/images/MarsHilUniversityLogo-Virtical.png}}

\institute{Mars Hill University}

\makebeamertitle 

%%True false command
\newcommand{\T}{\textcolor{green}{\contour{black}{T}}}
\newcommand{\F}{\textcolor{red}{\contour{black}{F}}}
\newcommand{\uT}[2]{\uncover<#1-#2>{\T}}
\newcommand{\uF}[2]{\uncover<#1-#2>{\F}}

%tkiz ball item 
\newcommand*\circled[1]{\tikz[baseline=(char.base)]
	{\node[circle,ball color=green!50!black!100, shade,   color=white,inner sep=1.2pt] (char) {\tiny #1};
	}
}
\begin{frame}{Variations of a Conditional}

\vspace{-.25cm}
\begin{block}{Variations of $p\rightarrow q$}

\begin{itemize}
\item [] \center%
\begin{tabular}{|c|c|c|}
\hline 
\rowcolor{blue!50}\textbf{Name} & \textbf{Symbol} & \textbf{Read as...}\tabularnewline
\hline 
\only<2->{\emph{Conditional}} & \only<2->{$p\rightarrow q$} & \only<2->{If $p$ then $q$}\tabularnewline
\hline 
\only<3->{\emph{Converse }(of $p\rightarrow q$)} & \only<4-4>{$p\leftarrow q$}\only<5->{$q\rightarrow p$} & \only<6->{if $q$ then $p$}\tabularnewline
\hline 
\only<7->{\emph{Inverse }(of $p\rightarrow q$)} & \uncover<8->{$\sim p\rightarrow\sim q$} & \uncover<9->{If not $p$ then not $q$}\tabularnewline
\hline 
\uncover<10->{\emph{Contrapositive }(of $p\rightarrow q$)} & \only<11->{$\sim q\rightarrow\sim p$} & \uncover<12->{If not $q$ then not $p$}\tabularnewline
\hline 
\end{tabular}
\end{itemize}
\end{block}
\vspace{-.5cm}
\begin{itemize}
\item [] \only<13-31>{%
\begin{tabular}{cc|c|c|c|c|c|c}
 & \multicolumn{1}{c}{} & \multicolumn{1}{c}{} & \multicolumn{1}{c}{} & \multicolumn{1}{c}{\textbf{Cond.}} & \multicolumn{1}{c}{\textbf{Contra.}} & \multicolumn{1}{c}{\textbf{Conv.}} & \textbf{Inv.}\tabularnewline
\textbf{$p$} & \textbf{$q$} & $\sim p$ & $\sim q$ & $p\rightarrow q$ & $\sim q\rightarrow\sim p$ & $q\rightarrow p$ & $\sim p\rightarrow\sim q$\tabularnewline
\hline 
\uT{13}{} & \uT{13}{} & \uF{14}{} & \uF{15}{} & \uT{16}{} & \uT{20}{} & \uT{24}{} & \uT{28}{}\tabularnewline
\uT{13}{} & \uF{13}{} & \uF{14}{} & \uT{15}{} & \uF{17}{} & \uF{21}{} & \uT{25}{} & \uT{29}{}\tabularnewline
\uF{13}{} & \uT{13}{} & \uT{14}{} & \uF{15}{} & \uT{18}{} & \uT{22}{} & \uF{26}{} & \uF{30}{}\tabularnewline
\uF{13}{} & \uF{13}{} & \uT{14}{} & \uT{15}{} & \uT{19}{} & \uT{23}{} & \uT{27}{} & \uT{31}{}\tabularnewline
\end{tabular}}\only<32-32>{%
\begin{tabular}{cc|c|c|c|c|c|c}
 & \multicolumn{1}{c}{} & \multicolumn{1}{c}{} & \multicolumn{1}{c}{} & \multicolumn{1}{c}{\cellcolor{blue!35}\textbf{Cond.}} & \multicolumn{1}{c}{\cellcolor{blue!35}\textbf{Contra.}} & \multicolumn{1}{c}{\textbf{Conv.}} & \textbf{Inv.}\tabularnewline
\textbf{$p$} & \textbf{$q$} & $\sim p$ & $\sim q$ & \cellcolor{blue!35}$p\rightarrow q$ & \cellcolor{blue!35}$\sim q\rightarrow\sim p$ & $q\rightarrow p$ & $\sim p\rightarrow\sim q$\tabularnewline
\hline 
\uT{13}{} & \uT{13}{} & \uF{14}{} & \uF{15}{} & \cellcolor{blue!35}\uT{16}{} & \cellcolor{blue!35}\uT{20}{} & \uT{24}{} & \uT{28}{}\tabularnewline
\uT{13}{} & \uF{13}{} & \uF{14}{} & \uT{15}{} & \cellcolor{blue!35}\uF{17}{} & \cellcolor{blue!35}\uF{21}{} & \uT{25}{} & \uT{29}{}\tabularnewline
\uF{13}{} & \uT{13}{} & \uT{14}{} & \uF{15}{} & \cellcolor{blue!35}\uT{18}{} & \cellcolor{blue!35}\uT{22}{} & \uF{26}{} & \uF{30}{}\tabularnewline
\uF{13}{} & \uF{13}{} & \uT{14}{} & \uT{15}{} & \cellcolor{blue!35}\uT{19}{} & \cellcolor{blue!35}\uT{23}{} & \uT{27}{} & \uT{31}{}\tabularnewline
\end{tabular}}\only<33->{%
\begin{tabular}{cc|c|c|c|c|c|c}
 & \multicolumn{1}{c}{} & \multicolumn{1}{c}{} & \multicolumn{1}{c}{} & \multicolumn{1}{c}{\cellcolor{blue!35}\textbf{Cond.}} & \multicolumn{1}{c}{\cellcolor{blue!35}\textbf{Contra.}} & \multicolumn{1}{c}{\cellcolor{red!35}\textbf{Conv.}} & \cellcolor{red!35}\textbf{Inv.}\tabularnewline
\textbf{$p$} & \textbf{$q$} & $\sim p$ & $\sim q$ & \cellcolor{blue!35}$p\rightarrow q$ & \cellcolor{blue!35}$\sim q\rightarrow\sim p$ & \cellcolor{red!35}$q\rightarrow p$ & \cellcolor{red!35}$\sim p\rightarrow\sim q$\tabularnewline
\hline 
\uT{13}{} & \uT{13}{} & \uF{14}{} & \uF{15}{} & \cellcolor{blue!35}\uT{16}{} & \cellcolor{blue!35}\uT{20}{} & \cellcolor{red!35}\uT{24}{} & \cellcolor{red!35}\uT{28}{}\tabularnewline
\uT{13}{} & \uF{13}{} & \uF{14}{} & \uT{15}{} & \cellcolor{blue!35}\uF{17}{} & \cellcolor{blue!35}\uF{21}{} & \cellcolor{red!35}\uT{25}{} & \cellcolor{red!35}\uT{29}{}\tabularnewline
\uF{13}{} & \uT{13}{} & \uT{14}{} & \uF{15}{} & \cellcolor{blue!35}\uT{18}{} & \cellcolor{blue!35}\uT{22}{} & \cellcolor{red!35}\uF{26}{} & \cellcolor{red!35}\uF{30}{}\tabularnewline
\uF{13}{} & \uF{13}{} & \uT{14}{} & \uT{15}{} & \cellcolor{blue!35}\uT{19}{} & \cellcolor{blue!35}\uT{23}{} & \cellcolor{red!35}\uT{27}{} & \cellcolor{red!35}\uT{31}{}\tabularnewline
\end{tabular}}
\item <13->Some are equivalent, and some are \textbf{not}.
\end{itemize}
\end{frame}

\begin{frame}{Variations of a Conditional}

\vspace{-.25cm}
\begin{block}{Variations of $p\rightarrow q$}

\begin{itemize}
\item [] \center%
\begin{tabular}{|c|c|c|}
\hline 
\rowcolor{blue!50}\textbf{Name} & \textbf{Symbol} & \textbf{Read as...}\tabularnewline
\hline 
\emph{Conditional} & $p\rightarrow q$ & If $p$ then $q$\tabularnewline
\hline 
\emph{Converse }(of $p\rightarrow q$) & $q\rightarrow p$ & if $q$ then $p$\tabularnewline
\hline 
\emph{Inverse }(of $p\rightarrow q$) & $\sim p\rightarrow\sim q$ & If not $p$ then not $q$\tabularnewline
\hline 
\emph{Contrapositive }(of $p\rightarrow q$) & $\sim q\rightarrow\sim p$ & If not $q$ then not $p$\tabularnewline
\hline 
\end{tabular}
\end{itemize}
\end{block}
\vspace{0cm}

\ex{2}{8}{If you pass this course, then you do graduate.
\begin{itemize}
\item []<3->\textbf{Converse:} \uncover<4->{''If you graduate then you
pass the course.''}
\item []<5->\textbf{Inverse: }\uncover<6->{''If you do not pass this
course then you do not graduate.''}
\item []<7->\textbf{Contrapositive:} \uncover<8->{''If you do not graduate
then you do not pass the course.''}
\end{itemize}
}

\ex{10}{17}{You are an orthodontist \only<10-10>{\emph{only if}}\only<11->{$\underset{\longrightarrow}{\underbrace{\mbox{\emph{only\ensuremath{} if}}}}$}\emph{
}you are a dentist.
\begin{itemize}
\item []<12->\textbf{Converse:} \uncover<13->{''If you are a dentist
then you are an orthodontist.''}
\item []<14->\textbf{Inverse:} \uncover<15->{''If you are not an orthodontist
then you are not a dentist.''}
\item []<16->\textbf{Contrapositive:} \uncover<17->{''If you are not
a dentist then you are not an orthodontist.''}
\end{itemize}
}

\ex{18}{}{

The TV does not work if the electricity is turned off.
\begin{itemize}
\item []<19->''If the TV does work then the electricity is turned on''.
\uncover<20->{\textbf{Contrapostive}}
\item []<19->''If the electricity is turned on then the TV does work''.
\uncover<21->{\textbf{Inverse}}
\item []<19->''If the TV does not work then the electricity is turned
off''. \uncover<22->{\textbf{Converse}}
\end{itemize}
}
\end{frame}

\begin{frame}{Writing a Contrapositive}

Writing a contrapositive of a given statement in terms of:
\begin{itemize}
\item <1->\emph{sufficient condition}
\item <16->\emph{necessary condition}
\end{itemize}
\ex{1}{}{Being alive is sufficient for running for office.\footnotesize \center

\begin{tabular}{ccc}
\only<2->{If \Ovalbox{a person is alive}} & \only<3->{then} & \only<4->{\Ovalbox{the person can run for office.}} \tabularnewline
\only<5->{\textcolor{red}{$p$: premise}} & \only<6->{\textcolor{red}{$\longrightarrow$}} & \only<7->{\textcolor{red}{$q$: conclusion}}\tabularnewline
\only<10->{If \Ovalbox{she can \emph{not} run for office} & \only<11->{then} & \only<12->{\Ovalbox{she is \emph{not} alive}.}} \tabularnewline
\only<9->{\textcolor{red}{$\sim q$: premise}} & \only<9->{\textcolor{red}{$\longrightarrow$}} & \only<9->{\textcolor{red}{$\sim p$: conclusion}}\tabularnewline
\only<13-16>{\Ovalbox{\emph{Not} running for office}} \only<17->{\Ovalbox{\emph{Not} being alive}} & \only<14-16>{is sufficient} \only<17->{is necessary} & \only<15-16>{\Ovalbox{for \emph{not} being alive.}} \only<17->{\Ovalbox{for \emph{not} running for office.}}\tabularnewline
\end{tabular}

}
\end{frame}

\begin{frame}{Biconditional }

\ex{1}{}{
\begin{itemize}
\item []<1->You get a refund if you have a receipt.
\item []<5->You get a refund only if you have a receipt.
\item []<7->You get a refund \emph{if and only if} you have a receipt.
\end{itemize}
}
\begin{columns}


\column{5cm}
\begin{itemize}
\item []<2->$p$: You get a refund.
\item []<3->$q:$ You have a receipt
\end{itemize}

\column{6cm}
\begin{itemize}
\item []<4->$q\rightarrow p$
\item []<6->$p\rightarrow q$
\item []<8->\Ovalbox{$p\leftrightarrow q$}
\end{itemize}
\end{columns}
\vspace{0cm}
\begin{itemize}
\item []<9->%
\begin{tabular}{c|c|c|c|c}
\textbf{$p$} & \textbf{$q$} & $p\rightarrow q$ & $q\rightarrow p$ & $(p\rightarrow q)\wedge(q\rightarrow q)$\tabularnewline
\hline 
\uT{9}{} & \uT{9}{} & \uT{10}{} & \uT{14}{} & \uT{18}{}\tabularnewline
\uT{9}{} & \uF{9}{} & \uF{11}{} & \uT{15}{} & \uF{19}{}\tabularnewline
\uF{9}{} & \uT{9}{} & \uT{12}{} & \uF{16}{} & \uF{20}{}\tabularnewline
\uF{9}{} & \uF{9}{} & \uT{13}{} & \uT{17}{} & \uT{21}{}\tabularnewline
\end{tabular}
\end{itemize}
\end{frame}

\begin{frame}{Problems}

\prob{1}{1}{Express the given biconditional statement as the conjuction of two conditionals}{A polygon is a triangle if and only if the polygon has three sides.}

\prob{2}{2}{Translate the two statements into symbolic form and use truth tables to determine whether the statements are equivalent.}{
\begin{itemize}
\item [1.] I watch TV only if the program is educational.
\item [2.] I do not watch TV if the program is not educational.
\end{itemize}
}

\prob{3}{3}{Express the contrapositive of the conditional in terms of a sufficient and necessary condition.}{Knowing
Morse code is necessary for operating a telegraph.}
\end{frame}

\end{document}
